\subsection{\added{Excluded Python features}}
\label{sec:excluded-features}

\begin{AddedBlock}
A syntactically valid Python program may use an \emph{excluded Python feature}: a syntactic form, module or module
member which \PurePy does not include. Each excluded feature is either \emph{excluded by design}, because it conflicts
with the design goals of \PurePy, such as immutability, or \emph{excluded from this version}, meaning it is likely to
appear in a future version of the language. \figref{excluded-by-design} lists the features excluded by design; any
other excluded feature is excluded from this version, and \figref{excluded-from-version} in the informative
\annexref{exclusions} lists the common ones.

A program may also be rejected because it uses an included feature in a way which the well-formedness rules do not
allow, such as a non-Boolean test. Such a rejection is by design, with the exception of the uses listed in
\figref{excluded-uses}, which are excluded from this version.
\end{AddedBlock}

\begin{figure}
\centering
\footnotesize
\setlength{\tabcolsep}{4pt}
\newcommand*{\designgroup}[1]{\rowcolor{grammarheadcolor}\multicolumn{3}{l}{\itshape #1} \\}
\begin{AddedBlock}
\begin{tabular}{>{\raggedright\arraybackslash}p{0.42\textwidth}>{\raggedright\arraybackslash}p{0.34\textwidth}>{\raggedright\arraybackslash}p{0.14\textwidth}}
\arrayrulecolor{gray!60}\hline
\rowcolor{grammarheadcolor}Feature & Python & Rejected by \\
\designgroup{Mutable state}
\pyinline{global} and \pyinline{nonlocal} & Assigns to variable of enclosing scope & Parser \\\rowgap
\pyinline{del} & Removes binding & Parser \\\rowgap
Augmented assignment \pyinline{+=}, etc. & Updates in place if type supports it & Parser \\\rowgap
Walrus operator \pyinline{:=} & Assigns within expression & Parser \\\rowgap
Item assignment \pyinline{xs[0] = v} & Updates container in place & Parser \\\rowgap
Attribute as assignment or declaration target, \pyinline{p.x = v} or \pyinline{p.x: int} & Updates object in place & Parser \\\rowgap
\pyinline{is} and \pyinline{is not}, other than against \pyinline{None}, \pyinline{True} or \pyinline{False} & Compares objects by identity & Parser \\\rowgap
Built-in \pyinline{id} & Gives identity of object & Type checker \\\rowgap
Mutating methods of built-in types: \pyinline{append}, \pyinline{extend}, \pyinline{insert}, \pyinline{pop},
\pyinline{remove}, \pyinline{clear}, \pyinline{sort} and \pyinline{reverse} of \pyinline{list}, and \pyinline{clear},
\pyinline{pop}, \pyinline{popitem}, \pyinline{setdefault} and \pyinline{update} of \pyinline{dict} & Update value in
place & Type checker \\\rowgap
\designgroup{Control flow and effects}
\pyinline{for} and \pyinline{while} loops, \pyinline{break} and \pyinline{continue} & Iterates by mutating loop state & Parser \\\rowgap
\pyinline{try}, \pyinline{except}, \pyinline{finally} and \pyinline{raise} & Raises and handles exceptions & Parser \\\rowgap
\pyinline{with} & Enters and exits context manager & Parser \\\rowgap
\pyinline{yield}, \pyinline{yield from} and generator expressions & Creates generator, consumed by mutation & Parser \\\rowgap
\pyinline{async} and \pyinline{await} & Defines and awaits coroutines & Parser \\\rowgap
Built-ins \pyinline{iter} and \pyinline{next} & Creates and advances iterator & Type checker \\\rowgap
Built-ins \pyinline{input} and \pyinline{open} & Reads input or opens file & Type checker \\\rowgap
\designgroup{Load-dependent meaning}
Import inside function body or branch, or after a statement & Dependencies of module known only as its statements run & Parser \\\rowgap
Wildcard import \pyinline{from q import *} & Introduces names of \pyinline{q} implicitly & Parser \\\rowgap
Built-ins \pyinline{eval}, \pyinline{exec}, \pyinline{getattr}, \pyinline{setattr}, \pyinline{delattr},
\pyinline{globals}, \pyinline{locals}, \pyinline{vars} and \pyinline{__import__} & Resolves or binds names at run time & Type checker \\\rowgap
\designgroup{Implicit coercions}
Dictionary type with key type other than \pyinline{str}, or dictionary pattern with key other than string literal & Accepts any hashable key & Parser \\\rowgap
\designgroup{Type annotations}
Parameter or return type without annotation & Accepts definition without annotations & Parser \\\rowgap
Annotation which is not a type expression under the Python typing specification, such as \pyinline{list[int, int]} or \pyinline{Literal[1.5]} & Ignores annotation & Parser \\\rowgap
\pyinline{Any} and \pyinline{cast} of \pyinline{typing} & Disables static checking of value & Type checker
\end{tabular}
\end{AddedBlock}
\caption{\added{Features outside \PurePy excluded by design}}
\label{fig:excluded-by-design}
\end{figure}


\begin{figure}
\centering
\footnotesize
\setlength{\tabcolsep}{4pt}
\begin{AddedBlock}
\begin{tabular}{>{\raggedright\arraybackslash}p{0.42\textwidth}>{\raggedright\arraybackslash}p{0.48\textwidth}}
\arrayrulecolor{gray!60}\hline
\rowcolor{grammarheadcolor}Use & Python \\
\pyinline{print} inside function or as first-class value & Print, or denote built-in function \\\rowgap
Second argument of \pyinline{round} & Round to that number of digits \\\rowgap
Argument of \pyinline{sys.exit} other than integer or Boolean, or no argument & Exit with status 0 for \pyinline{None} or no
argument; otherwise print argument and exit with status 1
\end{tabular}
\end{AddedBlock}
\caption{\added{Uses of predefined functions excluded from this version}}
\label{fig:excluded-uses}
\end{figure}

